% ===========================================================
\section{Arquitectura de Datos y Procesamiento Analítico}
% ===========================================================

\subsection{Motivación y Contexto}

El modelo operacional de \textbf{Bouquet} está diseñado para un entorno OLTP
(\textit{Online Transaction Processing}), caracterizado por alta concurrencia,
escrituras frecuentes, integridad referencial estricta y consultas de baja latencia.
Este diseño resulta adecuado para la operación diaria del restaurante, incluyendo
apertura de mesas, activación de sesiones virtuales, captura de órdenes, cambios de
estado en cocina, liquidación interna y cierre operativo de mesas.

Sin embargo, las consultas analíticas históricas requieren operaciones diferentes a las
consultas transaccionales. Para obtener métricas de consumo, productos más solicitados,
tiempos de preparación, comportamiento por franja horaria o estimaciones exploratorias
de demanda, es necesario realizar agregaciones sobre múltiples dominios: órdenes,
ítems, sesiones de mesa, comensales, estados de cocina, aportaciones internas,
ajustes autorizados, productos, categorías, estaciones y sucursales.

Ejecutar estas consultas directamente sobre la base transaccional cada vez que se carga
el dashboard podría afectar el rendimiento percibido por los usuarios operativos. Por
ello, Bouquet incorpora una capa de procesamiento analítico separada, basada en
Apache Spark, que extrae datos operacionales de forma controlada, los transforma bajo
una arquitectura Medallion y almacena resultados precalculados en tablas analíticas
Gold dentro de PostgreSQL.

De esta manera, el dashboard administrativo consulta tablas preparadas para análisis,
en lugar de recalcular métricas pesadas sobre el modelo OLTP. Esta separación sigue el
principio de CQRS (\textit{Command Query Responsibility Segregation}) aplicado al diseño
de datos: las operaciones transaccionales se atienden mediante el modelo operacional,
mientras que las consultas históricas y agregadas se atienden mediante estructuras
analíticas derivadas.

\subsection{Visión General de la Arquitectura}

La arquitectura analítica de Bouquet sigue el patrón \textbf{Medallion Architecture},
el cual organiza el procesamiento de datos en tres capas progresivas: Bronze, Silver
y Gold. Cada capa cumple una función específica dentro del refinamiento del dato, desde
la extracción cruda hasta la generación de métricas listas para consumo por el dashboard.

\begin{figure}[H]
  \centering
  \includegraphics[width=0.95\textwidth,height=0.4\textheight,keepaspectratio]{../_assets/spark/rendered/Spark-Ejecutivo.pdf}
  \caption{Arquitectura analítica de Bouquet --- Vista ejecutiva. El flujo inicia en las tablas transaccionales de PostgreSQL, continúa con el procesamiento batch en Apache Spark y termina en tablas Gold consultadas por el dashboard administrativo.}
  \label{fig:spark-ejecutivo}
\end{figure}

El flujo general opera de la siguiente forma:

\begin{enumerate}[leftmargin=1.5em]
  \item El dashboard administrativo solicita métricas históricas o dispara una actualización
  analítica autorizada.
  \item Next.js envía la solicitud al \textit{Analytics Orchestrator}, desplegado en el
  servidor VPS.
  \item El orquestador crea un registro en \texttt{AnalyticsJobRun}, indicando tipo de job,
  restaurante procesado, estado inicial y ventana de datos.
  \item El orquestador ejecuta los jobs de Spark mediante \texttt{spark-submit}.
  \item Spark extrae datos transaccionales hacia Bronze, construye datos limpios en Silver
  y genera agregados Gold.
  \item Las tablas Gold se escriben en PostgreSQL junto con el \texttt{jobRunId} que permite
  rastrear su origen.
  \item El dashboard consulta las tablas Gold mediante Prisma y renderiza las métricas
  sin ejecutar agregaciones pesadas sobre las tablas OLTP.
\end{enumerate}

Este diseño permite que el procesamiento analítico se ejecute de forma programada o bajo
solicitud autorizada, sin bloquear la experiencia operativa de meseros, cocina, comensales
o administradores.

\subsection{Arquitectura Medallion}

\subsubsection{Principio General}

La arquitectura Medallion organiza los datos en tres capas de refinamiento progresivo:

\begin{itemize}[leftmargin=1.5em]
  \item \textbf{Bronze} (\textit{Raw Data}): contiene copias crudas de los datos
  transaccionales. Conserva identificadores, marcas temporales y relaciones mínimas
  necesarias para reprocesar información sin volver a consultar todo el modelo operacional.

  \item \textbf{Silver} (\textit{Cleaned and Enriched Data}): contiene datos limpios,
  normalizados y enriquecidos mediante relaciones entre dominios. En esta capa se resuelven
  joins, se estandarizan estados, se normalizan fechas y se preparan campos para agregación.

  \item \textbf{Gold} (\textit{Business Ready Data}): contiene métricas agregadas listas
  para consumo por el dashboard. Estas tablas reducen la carga de consulta sobre el modelo
  OLTP y permiten mostrar indicadores de negocio de forma eficiente.
\end{itemize}

Las capas Bronze, Silver y Gold se almacenan como archivos Apache Parquet en el VPS, bajo
una estructura particionada por fecha de ejecución y tipo de job. Las tablas Gold, además,
se escriben de vuelta en PostgreSQL para que Next.js pueda consultarlas mediante Prisma.

\begin{verbatim}
/data/
  bronze/YYYY-MM-DD/
  silver/YYYY-MM-DD/
  gold/YYYY-MM-DD/
\end{verbatim}

La partición por fecha permite mantener historial de ejecuciones, facilita la depuración
de resultados y habilita reprocesamiento por ventanas específicas.

\subsubsection{Capa Bronze --- Ingesta Cruda}

La capa Bronze contiene datos crudos extraídos desde PostgreSQL mediante JDBC. Su objetivo
es capturar una fotografía de las entidades transaccionales necesarias para el análisis,
sin aplicar transformaciones de negocio.

\begin{longtable}{p{6cm}p{10cm}}
\caption{Fuentes de datos para la capa Bronze} \label{tab:bronze_sources} \\
\hline
\textbf{Archivo Parquet} & \textbf{Tabla(s) fuente} \\
\hline
\endfirsthead
\caption[]{TABLA \thetable{} (Continuación)} \\
\hline
\textbf{Archivo Parquet} & \textbf{Tabla(s) fuente} \\
\hline
\endhead
\hline
\endfoot
\texttt{bronze\_orders} &
\texttt{RestaurantOrder} \\
\hline
\texttt{bronze\_order\_items} &
\texttt{OrderItem}, \texttt{OrderItemModifier} \\
\hline
\texttt{bronze\_status\_events} &
\texttt{OrderItemStatusEvent} \\
\hline
\texttt{bronze\_settlements} &
\texttt{Settlement}, \texttt{SettlementContribution}, \texttt{SettlementAllocation},
\texttt{AppliedAdjustment} \\
\hline
\texttt{bronze\_menu} &
\texttt{RestaurantMenuItem}, \texttt{RestaurantCategory}, \texttt{ModifierGroup},
\texttt{ModifierOption} \\
\hline
\texttt{bronze\_sessions} &
\texttt{DiningSession}, \texttt{DiningSessionTable}, \texttt{DiningTable}, \texttt{Guest} \\
\hline
\texttt{bronze\_catalog} &
\texttt{Chain}, \texttt{Zone}, \texttt{Restaurant}, \texttt{Station} \\
\hline
\end{longtable}

La extracción puede realizarse de forma incremental utilizando las marcas temporales
\texttt{createdAt}, \texttt{updatedAt} o campos específicos de eventos, según la tabla.
La ventana procesada se registra en \texttt{AnalyticsJobRun.dataWindowStart} y
\texttt{AnalyticsJobRun.dataWindowEnd}, lo que permite reprocesar periodos específicos
sin recalcular todo el histórico.

\subsubsection{Capa Silver --- Limpieza y Enriquecimiento}

La capa Silver transforma los datos crudos de Bronze en datasets limpios y enriquecidos.
Esta capa funciona como el punto intermedio entre el modelo transaccional normalizado y
las métricas de negocio.

Las transformaciones principales son:

\begin{enumerate}[leftmargin=1.5em]
  \item \textbf{Resolución de relaciones entre dominios:} se integran órdenes, ítems,
  sesiones, productos, categorías, estaciones, restaurantes y liquidaciones internas.

  \item \textbf{Normalización de estados:} los estados operativos, como
  \texttt{PENDIENTE}, \texttt{EN\_PREPARACION}, \texttt{LISTA}, \texttt{ENTREGADA},
  \texttt{ABIERTA}, \texttt{PARCIAL} y \texttt{LIQUIDADA}, se transforman en etiquetas
  homogéneas para análisis.

  \item \textbf{Gestión de fechas y zona horaria:} las marcas temporales se estandarizan
  a la zona \texttt{America/Mexico\_City} para evitar errores en cortes diarios,
  horarios pico y ventanas de operación.

  \item \textbf{Conservación de montos en centavos:} los campos monetarios se mantienen
  como enteros en centavos para cálculo. La conversión a pesos se realiza únicamente en
  campos derivados de visualización.

  \item \textbf{Preparación de eventos de cocina:} a partir de
  \texttt{OrderItemStatusEvent} se calculan tiempos entre estados, como inicio de
  preparación, finalización y entrega.

  \item \textbf{Preparación de liquidación interna:} se consolidan aportaciones,
  asignaciones y ajustes autorizados para calcular saldos, montos netos y comportamiento
  de cierre.
\end{enumerate}

La capa Silver evita que las tablas Gold tengan que resolver joins complejos sobre el
modelo OLTP y permite reutilizar datasets limpios para diferentes métricas.

\subsubsection{Capa Gold --- Agregados de Negocio}

La capa Gold contiene las tablas analíticas listas para el dashboard. A diferencia de
Bronze y Silver, que viven principalmente como Parquet en el VPS, los resultados Gold se
escriben también en PostgreSQL. Esto permite que la aplicación Next.js consulte métricas
mediante Prisma, sin depender de Spark durante la carga del dashboard.

\begin{figure}[H]
  \centering
  \includegraphics[width=0.95\textwidth,height=0.85\textheight,keepaspectratio]{../_assets/spark/rendered/Spark-Tecnico.pdf}
  \caption{Pipeline técnico de Spark --- Vista detallada. Se muestran las capas Bronze, Silver y Gold, así como la escritura de métricas hacia tablas analíticas en PostgreSQL.}
  \label{fig:spark-tecnico}
\end{figure}

Las tablas Gold consideradas son:

\begin{itemize}[leftmargin=1.5em]
  \item \texttt{analytic\_sales\_daily}
  \item \texttt{analytic\_item\_velocity}
  \item \texttt{analytic\_service\_times}
  \item \texttt{analytic\_demand\_estimate}
\end{itemize}

Cada registro Gold conserva \texttt{jobRunId}, lo que permite rastrear qué ejecución del
pipeline generó la métrica.

\subsection{Jobs del Pipeline}

El pipeline se compone de cinco jobs secuenciales coordinados por el
\textit{Analytics Orchestrator}. Cada job actualiza su estado en \texttt{AnalyticsJobRun}
o en los metadatos asociados de ejecución.

\begin{enumerate}[leftmargin=1.5em]
  \item \textbf{Job 1 --- Extract Bronze:} extrae datos transaccionales desde PostgreSQL
  hacia archivos Parquet en la capa Bronze.

  \item \textbf{Job 2 --- Build Silver:} limpia, normaliza y enriquece datos Bronze,
  resolviendo relaciones entre órdenes, ítems, sesiones, liquidaciones, menú, restaurantes
  y estaciones.

  \item \textbf{Job 3 --- Aggregate Gold Metrics:} genera métricas de consumo diario,
  velocidad de productos y tiempos de servicio.

  \item \textbf{Job 4 --- Demand Estimate:} calcula estimaciones exploratorias de demanda
  mediante métodos simples e interpretables sobre históricos.

  \item \textbf{Job 5 --- Write Back:} escribe tablas Gold en PostgreSQL y actualiza el
  estado final del job analítico.
\end{enumerate}

\subsubsection{Job 1 --- Extract Bronze}

Este job consulta las tablas transaccionales mediante JDBC y genera archivos Parquet
crudos. La extracción se limita a la ventana definida por \texttt{dataWindowStart} y
\texttt{dataWindowEnd}, cuando dicha ventana está disponible.

Su propósito es reducir el número de consultas directas sobre PostgreSQL y permitir que
las transformaciones posteriores se realicen sobre archivos locales del VPS.

\subsubsection{Job 2 --- Build Silver}

Este job lee los archivos Bronze y produce datasets Silver limpios y enriquecidos. Entre
sus responsabilidades se encuentran:

\begin{itemize}[leftmargin=1.5em]
  \item unir órdenes con ítems, comensales, sesiones y restaurantes;
  \item unir ítems con productos, categorías y estaciones;
  \item integrar eventos de estado para reconstruir tiempos de cocina;
  \item integrar liquidaciones, aportaciones, asignaciones y ajustes;
  \item normalizar estados, fechas y montos;
  \item descartar o marcar registros incompletos que no puedan utilizarse en agregados.
\end{itemize}

\subsubsection{Job 3 --- Aggregate Gold Metrics}

Este job produce las tablas \texttt{analytic\_sales\_daily},
\texttt{analytic\_item\_velocity} y \texttt{analytic\_service\_times}. Sus agregaciones
se basan en datos Silver enriquecidos, evitando que el dashboard ejecute joins históricos
sobre la base OLTP.

Las métricas generadas incluyen:

\begin{itemize}[leftmargin=1.5em]
  \item número de órdenes registradas por fecha operativa;
  \item número de ítems solicitados;
  \item monto bruto de consumo registrado;
  \item monto reducido por descuentos o cortesías autorizadas;
  \item monto neto de consumo registrado;
  \item productos más solicitados;
  \item tiempos promedio de preparación;
  \item percentiles de preparación por estación;
  \item ítems retrasados según umbral operativo.
\end{itemize}

\subsubsection{Job 4 --- Demand Estimate}

Este job produce la tabla \texttt{analytic\_demand\_estimate}. Su objetivo es generar
estimaciones exploratorias de demanda por producto y fecha. No se considera un sistema
predictivo avanzado, sino un mecanismo inicial para observar patrones históricos y apoyar
decisiones administrativas.

El método base propuesto es una media móvil ponderada con estacionalidad semanal. Para
cada producto se calcula el comportamiento histórico por día de la semana, asignando mayor
peso a los registros recientes. Este enfoque es simple, interpretable y adecuado para una
primera versión académica.

El campo \texttt{confidenceBps} se interpreta como un indicador opcional de confiabilidad
de la estimación. Puede calcularse a partir de la variabilidad histórica, por ejemplo
mediante una función basada en el coeficiente de variación:

\[
  \texttt{confidenceBps} = \frac{10000}{1 + CV}
\]

donde $CV$ representa el coeficiente de variación histórico. Al utilizar basis points,
un valor de $10000$ representa una confiabilidad estimada del $100\%$.

En versiones futuras, este módulo podría reemplazarse por métodos estadísticos o de
aprendizaje automático más sofisticados sin modificar la interfaz de salida de la tabla
Gold.

\subsubsection{Job 5 --- Write Back}

El último job lee los DataFrames Gold y los escribe en tablas analíticas de PostgreSQL.
La escritura se realiza mediante una estrategia de reemplazo controlado por restaurante,
tabla y ventana de datos procesada.

Para evitar duplicados, cada tabla Gold define llaves lógicas, por ejemplo:

\begin{itemize}[leftmargin=1.5em]
  \item \texttt{restaurantId + businessDate} para \texttt{analytic\_sales\_daily};
  \item \texttt{restaurantId + menuItemId + businessDate} para \texttt{analytic\_item\_velocity};
  \item \texttt{restaurantId + stationId + businessDate} para \texttt{analytic\_service\_times};
  \item \texttt{restaurantId + menuItemId + estimateDate} para \texttt{analytic\_demand\_estimate}.
\end{itemize}

Cuando la tabla lo permite, se utiliza una operación tipo \textit{upsert} para actualizar
métricas existentes sin afectar datos fuera de la ventana procesada. Cada registro escrito
mantiene \texttt{jobRunId}, permitiendo rastrear su origen analítico.

\subsection{Seguridad y Control Operacional}

Dado que el procesamiento analítico puede consumir recursos considerables, el sistema
incorpora mecanismos de protección y control:

\begin{itemize}[leftmargin=1.5em]
  \item \textbf{Seguridad de API:} la comunicación entre Next.js y el
  \textit{Analytics Orchestrator} se protege mediante una \textit{API Key} compartida,
  almacenada en variables de entorno.

  \item \textbf{Control de concurrencia:} antes de iniciar un nuevo proceso, el
  orquestador consulta \texttt{AnalyticsJobRun}. Si existe un job en estado
  \texttt{QUEUED} o \texttt{RUNNING} para el mismo \texttt{restaurantId} y tipo de job,
  la solicitud se rechaza para evitar sobrecarga y condiciones de carrera.

  \item \textbf{Aislamiento de recursos:} Spark se configura con límites de memoria y
  núcleos de CPU para evitar consumo excesivo de recursos en el VPS.

  \item \textbf{Separación de ejecución:} el procesamiento analítico no se ejecuta dentro
  de Vercel ni dentro del proceso principal de la aplicación web; se ejecuta en el VPS
  mediante el orquestador.

  \item \textbf{Ventanas de procesamiento:} cada ejecución registra
  \texttt{dataWindowStart} y \texttt{dataWindowEnd}, permitiendo limitar el rango de datos
  procesado.
\end{itemize}

\subsection{Linaje de Datos y Transformaciones}

El linaje de datos permite explicar cómo se transforma la información desde el modelo
transaccional hasta las tablas Gold. Esto facilita validación, depuración y defensa de
métricas durante revisiones técnicas.

\subsubsection{Linaje de \texttt{analytic\_sales\_daily}}

\begin{figure}[H]
  \centering
  \includegraphics[width=0.95\textwidth,height=0.6\textheight,keepaspectratio]{../_assets/spark/rendered/Lineage-Sales.pdf}
  \caption{Linaje de \texttt{analytic\_sales\_daily}. El flujo consolida órdenes, ítems, sesiones, restaurantes y ajustes autorizados para generar métricas diarias de consumo.}
  \label{fig:lineage-sales}
\end{figure}

\begin{table}[H]
\centering
\caption{Diccionario de Métricas --- Sales Daily}
\begin{tabular}{p{4cm}p{5cm}p{7cm}}
\hline
\textbf{Métrica} & \textbf{Cálculo (Spark SQL)} & \textbf{Descripción} \\
\hline
\texttt{orderCount} &
\texttt{COUNT(DISTINCT orderId)} &
Número de órdenes registradas en la fecha operativa. \\
\hline
\texttt{itemCount} &
\texttt{SUM(quantity)} &
Número total de ítems registrados. \\
\hline
\texttt{grossConsumptionCents} &
\texttt{SUM(totalCents)} &
Monto bruto de consumo registrado. \\
\hline
\texttt{discountCents} &
\texttt{SUM(adjustmentAmountCents)} &
Monto reducido mediante descuentos o cortesías autorizadas. \\
\hline
\texttt{netConsumptionCents} &
\texttt{grossConsumptionCents - discountCents} &
Monto neto de consumo registrado después de ajustes. \\
\hline
\end{tabular}
\end{table}

\subsubsection{Linaje de \texttt{analytic\_item\_velocity}}

\begin{figure}[H]
  \centering
  \includegraphics[width=0.95\textwidth,height=0.6\textheight,keepaspectratio]{../_assets/spark/rendered/Lineage-Velocity.pdf}
  \caption{Linaje de \texttt{analytic\_item\_velocity}. Se enfoca en el comportamiento de consumo por producto, considerando cantidad solicitada y monto bruto asociado.}
  \label{fig:lineage-velocity}
\end{figure}

\begin{table}[H]
\centering
\caption{Diccionario de Métricas --- Item Velocity}
\begin{tabular}{p{4cm}p{6cm}p{6cm}}
\hline
\textbf{Métrica} & \textbf{Cálculo (Spark SQL)} & \textbf{Descripción} \\
\hline
\texttt{quantitySold} &
\texttt{SUM(quantity)} &
Cantidad solicitada del producto en el periodo. \\
\hline
\texttt{grossConsumptionCents} &
\texttt{SUM(totalCents)} &
Monto bruto asociado al producto. \\
\hline
\texttt{itemNameSnapshot} &
\texttt{LAST(itemNameSnapshot)} &
Nombre histórico usado para análisis aunque el catálogo cambie. \\
\hline
\end{tabular}
\end{table}

\subsubsection{Linaje de \texttt{analytic\_service\_times}}

\begin{figure}[H]
  \centering
  \includegraphics[width=0.95\textwidth,height=0.6\textheight,keepaspectratio]{../_assets/spark/rendered/Lineage-ServiceTimes.pdf}
  \caption{Linaje de \texttt{analytic\_service\_times}. Utiliza eventos de estado de \texttt{OrderItemStatusEvent} para calcular tiempos de preparación por estación a partir de las transiciones \texttt{PENDIENTE}, \texttt{EN\_PREPARACION} y \texttt{LISTA}.}
  \label{fig:lineage-service-times}
\end{figure}

\begin{table}[H]
\centering
\caption{Diccionario de Métricas --- Service Times}
\begin{tabular}{p{4cm}p{6cm}p{6cm}}
\hline
\textbf{Métrica} & \textbf{Cálculo (Spark SQL)} & \textbf{Descripción} \\
\hline
\texttt{avgPrepSeconds} &
\texttt{AVG(ts\_lista - ts\_preparacion)} &
Tiempo promedio de preparación. \\
\hline
\texttt{p50PrepSeconds} &
\texttt{PERCENTILE(duration, 0.50)} &
Mediana del tiempo de preparación. \\
\hline
\texttt{p90PrepSeconds} &
\texttt{PERCENTILE(duration, 0.90)} &
Percentil 90 del tiempo de preparación. \\
\hline
\texttt{delayedItemsCount} &
\texttt{COUNT(duration > threshold)} &
Ítems que superan el umbral operativo definido. \\
\hline
\end{tabular}
\end{table}

\subsubsection{Linaje de \texttt{analytic\_demand\_estimate}}

\begin{figure}[H]
  \centering
  \includegraphics[width=0.95\textwidth,height=0.6\textheight,keepaspectratio]{../_assets/spark/rendered/Lineage-DemandEstimate.pdf}
  \caption{Linaje de \texttt{analytic\_demand\_estimate}. Genera estimaciones exploratorias de demanda por producto y fecha a partir de históricos agregados por día de la semana.}
  \label{fig:lineage-demand-estimate}
\end{figure}

\begin{table}[H]
\centering
\caption{Diccionario de Métricas --- Demand Estimate}
\begin{tabular}{p{4cm}p{6cm}p{6cm}}
\hline
\textbf{Métrica} & \textbf{Cálculo (Spark SQL)} & \textbf{Descripción} \\
\hline
\texttt{expectedQuantity} &
\texttt{weighted\_avg(hist\_quantity)} &
Cantidad esperada estimada a partir de históricos. \\
\hline
\texttt{confidenceBps} &
\texttt{10000 / (1 + CV)} &
Indicador opcional de confiabilidad basado en variabilidad histórica. \\
\hline
\end{tabular}
\end{table}

\subsection{Analytics Orchestrator}

El \textbf{Analytics Orchestrator} es un servicio ligero implementado en \textbf{FastAPI}
que reside en el servidor VPS. Actúa como intermediario entre la aplicación Next.js y
el pipeline de Spark.

Sus responsabilidades son:

\begin{itemize}[leftmargin=1.5em]
  \item recibir solicitudes de ejecución desde la API de Next.js;
  \item generar un identificador único de job;
  \item crear y actualizar registros en \texttt{AnalyticsJobRun};
  \item ejecutar jobs de Spark mediante \texttt{spark-submit};
  \item rechazar ejecuciones concurrentes conflictivas;
  \item exponer el estado del job para el mecanismo de consulta del dashboard.
\end{itemize}

El servicio expone endpoints internos protegidos por API Key:

\begin{itemize}[leftmargin=1.5em]
  \item \texttt{POST /jobs}: inicia una ejecución del pipeline y retorna el identificador del job.
  \item \texttt{GET /jobs/\{jobId\}}: retorna el estado actual del job.
\end{itemize}

Estos endpoints no están diseñados para consumo directo por el cliente final; son invocados
por la capa de aplicación.

\subsection{Tablas Analíticas Gold en PostgreSQL}

Las tablas analíticas son el punto de integración entre Spark y el dashboard de Next.js.
Son escritas por el pipeline analítico y leídas por la capa de aplicación mediante Prisma.
Ninguna operación OLTP escribe directamente en estas tablas.

\begin{longtable}{p{6cm}p{10cm}}
\caption{Tablas analíticas Gold} \label{tab:gold_tables} \\
\hline
\textbf{Tabla} & \textbf{Contenido} \\
\hline
\endfirsthead
\caption[]{TABLA \thetable{} (Continuación)} \\
\hline
\textbf{Tabla} & \textbf{Contenido} \\
\hline
\endhead
\hline
\endfoot
\texttt{analytic\_sales\_daily} &
Métricas diarias de consumo por restaurante. Incluye órdenes registradas, ítems,
monto bruto de consumo, ajustes aplicados y monto neto de consumo. \\
\hline
\texttt{analytic\_item\_velocity} &
Velocidad de consumo por producto y día. Incluye cantidad solicitada, monto bruto
asociado y nombre histórico del producto. \\
\hline
\texttt{analytic\_service\_times} &
Tiempos de preparación por estación y día. Incluye promedio, percentil 50, percentil 90
e ítems retrasados según umbral operativo. \\
\hline
\texttt{analytic\_demand\_estimate} &
Estimaciones exploratorias de demanda por producto y fecha. Incluye cantidad esperada,
método utilizado e indicador opcional de confiabilidad. \\
\hline
\texttt{AnalyticsJobRun} &
Historial de ejecuciones del pipeline. Permite conocer estado, ventana procesada,
errores, metadatos y antigüedad del último análisis. \\
\hline
\end{longtable}

\subsection{Consideraciones de Escalabilidad}

La arquitectura analítica de Bouquet está diseñada para una primera versión académica y
operativa, con escala de restaurante individual o pequeña cadena. Sin embargo, incorpora
decisiones que facilitan su evolución futura:

\begin{itemize}[leftmargin=1.5em]
  \item \textbf{Particionamiento por fecha:} el almacenamiento Parquet particionado permite
  agregar nuevas ejecuciones sin reescribir datos históricos completos. En el futuro puede
  migrarse a almacenamiento distribuido como S3, GCS o HDFS sin cambiar la lógica central
  de los jobs.

  \item \textbf{Pipeline modular:} cada job es un componente independiente invocado por el
  orquestador. Esto permite agregar nuevos análisis, como inventario o segmentación de
  comportamiento, sin modificar todos los jobs existentes.

  \item \textbf{Estimación reemplazable:} el módulo de demanda utiliza un método simple e
  interpretable, pero puede reemplazarse por modelos estadísticos como ARIMA, Prophet o
  métodos de Spark MLlib manteniendo la misma tabla de salida.

  \item \textbf{Separación OLTP/OLAP:} el procesamiento analítico se ejecuta fuera del flujo
  operacional, evitando que el dashboard histórico afecte la captura de órdenes, estados
  de cocina o liquidaciones internas.

  \item \textbf{Diseño multi-tenant:} aunque la primera implementación puede procesar por
  \texttt{restaurantId}, el modelo organizacional permite extender análisis a nivel de zona
  o cadena mediante agregación sobre varias sucursales.
\end{itemize}